%========================== REFERÊNCIAS ==========================
\newpage
\section*{Referências}
\addcontentsline{toc}{section}{REFERÊNCIAS}

\begingroup
\singlespacing
\setlength{\parindent}{0pt}
\setlength{\parskip}{0.7\baselineskip}

BRASIL. \textbf{Lei n. 10.436, de 24 de abril de 2002}. Dispõe sobre a Língua
Brasileira de Sinais – Libras e dá outras providências. Diário Oficial da União:
Brasília, DF, 2002.

BRASIL. \textbf{Decreto n. 5.626, de 22 de dezembro de 2005}. Regulamenta a Lei
n. 10.436, de 24 de abril de 2002. Diário Oficial da União: Brasília, DF, 2005.

HOCHREITER, S.; SCHMIDHUBER, J. Long short-term memory. \textbf{Neural
Computation}, v. 9, n. 8, p. 1735--1780, 1997.

IBGE. [PENDENTE: referência completa da fonte do dado de 10 milhões de pessoas
com deficiência auditiva (pesquisa, ano e \textit{link})]

LANDIS, J. R.; KOCH, G. G. The measurement of observer agreement for categorical
data. \textbf{Biometrics}, v. 33, n. 1, p. 159--174, 1977.

MEDIAPIPE. \textbf{MediaPipe Hands Solution}. Google LLC. Disponível em:
\url{https://developers.google.com/mediapipe/solutions/vision/hand_landmarker}.
Acesso em: [PENDENTE: data].

MOREIRA, J. \textbf{Dynamic LIBRAS gesture recognition via CNN over spatiotemporal
matrix representation}. Curitiba: Universidade Tecnológica Federal do Paraná,
2026. Disponível em: arXiv:2603.25863.

PIGOU, L. et al. Beyond temporal pooling: recurrence and temporal convolutions
for gesture recognition in video. \textbf{International Journal of Computer
Vision}, v. 126, n. 2--4, p. 430--439, 2018.

[PENDENTE: referência completa do trabalho de reconhecimento de sílabas em Libras
para HRI (2026) citado na Tabela~\ref{tab:relacionados}]

QUADROS, R. M. de. \textbf{Libras}. São Paulo: Parábola, 2019. (Linguística para
o ensino superior).

QUADROS, R. M. de; KARNOPP, L. B. \textbf{Língua de sinais brasileira}: estudos
linguísticos. Porto Alegre: Artmed, 2004.

REGO, T. et al. Brazilian sign language recognition using deep learning based on
fast Fourier transform and kinematic features. \textbf{IEEE Access}, 2025.

REZENDE, T. M.; ALMEIDA, S. G. M.; GUIMARÃES, F. G. \textbf{Reconhecimento
automático de sinais da Libras}: desenvolvimento da base de dados MINDS-Libras e
modelos de redes neurais convolucionais. Tese (Doutorado em Engenharia Elétrica)
– Universidade Federal de Minas Gerais, Belo Horizonte, 2021.

RONCHETTI, F. et al. LSA64: an Argentinian sign language dataset. In: CONGRESO
ARGENTINO DE CIENCIAS DE LA COMPUTACIÓN, 22., 2016. \textbf{Actas} [...]. 2016.

SINCAN, O. M.; KELES, H. Y. AUTSL: a large scale multi-modal Turkish sign
language dataset and baseline methods. \textbf{IEEE Access}, v. 8, p.
181340--181355, 2020.

TENSORFLOW. \textbf{TensorFlow documentation}: Keras LSTM layer. Disponível em:
\url{https://www.tensorflow.org/api_docs/python/tf/keras/layers/LSTM}. Acesso em:
[PENDENTE: data].

ZHANG, F. et al. MediaPipe Hands: on-device real-time hand tracking. In: CVPR
WORKSHOP ON COMPUTER VISION FOR AUGMENTED AND VIRTUAL REALITY, 2020.
\textbf{Proceedings} [...]. 2020.
\endgroup
